\section{Instrument Validation}
\label{sec:instrument}

Before running the full budget we ran a 2{,}000-step pilot whose only purpose was to
check that the measurement apparatus can detect collapse when collapse is definitely
present. The \texttt{none\_nostopgrad} condition exists for exactly this: it has no EMA
split, no stop-gradient, and no regularizer, so nothing prevents the trivial constant
solution. If it does not register as collapsed, no other condition's number can be
trusted.

It collapsed immediately and completely. Two of the four diagnostics reported the
opposite.

\subsection{The control collapsed}

\begin{table}[h]
\centering
\small
\begin{tabular}{r r r r r r}
\toprule
\textbf{step} & \textbf{loss} & \textbf{feature std} & \textbf{cosine} &
\textbf{eff.\ rank} & \textbf{probe} \\
\midrule
0    & ---      & 0.623 & 0.333 & 2.71  & 0.375 \\
100  & 0.0021   & 0.024 & 0.999 & 9.06  & 0.412 \\
500  & 0.00009  & 0.005 & 1.000 & 23.58 & 0.422 \\
2000 & 0.00001  & 0.0019 & 1.000 & 30.29 & 0.409 \\
\bottomrule
\end{tabular}
\caption{\ours{} \texttt{none\_nostopgrad} over the pilot. Mean pairwise cosine reaches
1.000 and the prediction loss reaches $10^{-5}$: every input maps to essentially the same
vector. Effective rank and probe accuracy both \emph{rise}.}
\end{table}

By step 100 the mean pairwise cosine similarity is 0.999 and the per-feature standard
deviation has fallen by a factor of 25. By step 2000 the loss is $10^{-5}$. This is
textbook total collapse, and the condition was built to produce it.

\subsection{Two diagnostics moved the wrong way}

Effective rank rose from 2.71 to 30.29 and standardized probe accuracy rose from 0.375 to
0.409, both monotonically, as the representation died. \ours{} Both failures share one
cause: \textbf{scale invariance}.

\paragraph{Effective rank.} The participation ratio
$(\sum_i \lambda_i)^2 / \sum_i \lambda_i^2$ is a ratio of eigenvalue sums, so any overall
rescaling cancels. It describes the \emph{shape} of the spectrum and says nothing about
its magnitude. Once the signal variance falls below the numerical noise floor, the
covariance spectrum is dominated by floating-point noise, which is close to isotropic and
therefore has high participation ratio. The metric ends up measuring the shape of the
residue rather than the presence of signal. \interp{} We read this as effective rank being
a \emph{dimensional} collapse diagnostic that is being asked a question about
\emph{complete} collapse, which it is not equipped to answer alone.

\paragraph{The linear probe.} Following standard practice we standardized features before
fitting, so that the regularization strength means the same thing at every checkpoint as
embedding scale drifts. That choice is precisely what blinds the probe to scale collapse:
dividing by a per-feature standard deviation of $0.0019$ rescales the collapsed
representation's numerical residue back to unit variance, and that residue is still a
deterministic function of the input. A linear classifier reads it without difficulty.

We confirmed the mechanism directly on synthetic embeddings constructed as a large
constant plus a small input-dependent residue, with mean pairwise cosine similarity
$1.0000$:

\begin{table}[h]
\centering
\small
\begin{tabular}{l r}
\toprule
\textbf{probe variant} & \textbf{accuracy} \\
\midrule
with standardization & 1.000 \\
without standardization & 0.592 \\
chance (10 classes) & 0.100 \\
\bottomrule
\end{tabular}
\caption{\ours{} A totally collapsed representation is read perfectly by a standardized
linear probe. Standardization removes exactly the signal that collapse destroys.}
\end{table}

\subsection{Why this threatened the study, not just the control}

The condition under test, \texttt{sigreg\_nostopgrad}, was itself visibly degrading in the
pilot --- per-feature standard deviation fell from 0.623 to 0.295 and mean pairwise cosine
rose from 0.333 to 0.879 --- yet its standardized probe accuracy (0.379) read
\emph{higher} than that of the healthier-looking \texttt{sigreg\_stopgrad} (0.309).

More seriously, this report's Part 2 asks whether cheap diagnostics degrade \emph{before}
the probe does. If the probe does not degrade under collapse at all, that question is not
merely hard to answer, it is malformed as operationalised. \ours{}

\subsection{Changes made}

Three changes, made before any full-length run:

\begin{enumerate}
  \item \textbf{Log total variance} (the trace of the embedding covariance) as an explicit
        scale metric. This is the quantity that actually falls under collapse, and it is
        the necessary companion to effective rank.
  \item \textbf{Log probe accuracy twice}, standardized and unstandardized. The
        standardized variant remains comparable across checkpoints; the unstandardized
        variant is the one that registers scale collapse. A widening gap between them is
        itself a collapse signal.
  \item \textbf{Document that effective rank is not a standalone collapse detector},
        in its own docstring and wherever it is plotted. It is not wrong; it answers a
        different question than the one it is easily mistaken for.
\end{enumerate}

\established{} Both effective rank and standardized linear probing are widely used
collapse diagnostics, which is why this matters beyond one harness: the failure mode is a
property of the metrics, not of this implementation. \ours{} We are not aware of these two
being reported together as a joint blind spot, though we have not surveyed the literature
systematically enough to claim novelty, and the observation rests on a single-seed pilot
at one model scale.

\subsection{What the pilot does not establish}

\interp{} The pilot ran 2{,}000 steps on one seed. It is evidence that the apparatus was
broken and is now less broken; it is not evidence about the research question. In
particular, none of the between-condition orderings observed in the pilot should be read
as results --- they were measured with the instrument this section is about.

\subsection{A second defect: the regularizer was not regularizing}

The corrected instrument was then used for a second pilot, which produced a further
finding. Neither SIGReg condition exceeded random-initialisation probe accuracy (0.306 and
0.315 against 0.369 at step zero), while the EMA baseline reached 0.497. More tellingly,
\texttt{sigreg\_nostopgrad} slid into collapse over its final quarter --- mean pairwise
cosine 0.564 to 0.879, total variance 76 to 21 --- \emph{while SIGReg was numerically
dominant}, its weighted term running two to ten times the prediction loss.

A regularizer being paid that much and still losing is implausible, so we measured its
restoring force directly by scaling a batch toward a constant vector:

\begin{table}[h]
\centering
\small
\begin{tabular}{r r r}
\toprule
\textbf{residual noise} & \textbf{reference grad norm} & \textbf{our version} \\
\midrule
$10^{0}$  & 4.24 & $2.87\times10^{-4}$ \\
$10^{-2}$ & 3.64 & $2.53\times10^{-4}$ \\
$10^{-4}$ & 3.63 & $2.53\times10^{-6}$ \\
\bottomrule
\end{tabular}
\caption{\ours{} As collapse deepens, the reference implementation maintains its restoring
force while ours decays to nothing.}
\end{table}

\paragraph{Cause.} We centered the embeddings before the normality test. The reference does
not. A collapsed encoder emits some constant vector $c$; centering maps that exactly onto
the origin, where the characteristic function is flat ($\cos(0)=1$, derivative zero). The
loss saturated at a high value --- correctly reporting a severe problem --- while its
gradient vanished. Collapse had become a stationary point of the objective built to prevent
it, and the loss value and the loss gradient disagreed about how bad the situation was.

Centering additionally made a mean-shifted distribution invisible, since it removes exactly
the discrepancy a comparison against a \emph{standard} normal is meant to detect.

\paragraph{Corrections.} The module now matches the reference bit-for-bit on isotropic,
anisotropic, collapsed and mean-shifted inputs, verified by a test carrying a port of the
reference. Four further deviations were found and corrected: the integration grid excluded
$t=0$, the quadrature weights were normalized rather than trapezoidal, the batch-size
scaling of the Epps-Pulley statistic had been removed, and the defaults differed. Separately,
the loss weight was roughly twenty times the top of the reference's tuned range: LeJEPA
forms a convex combination and sweeps $\lambda \in \{0.01, 0.02, 0.05, 0.1\}$, equivalent to
$0.01$--$0.11$ in our additive parametrisation, where we had been using $1.0$.

\paragraph{Consequence.} \interp{} All SIGReg results from both pilots were produced by a
regularizer that was not regularizing, and none of them should be compared against anything.
This is also the strongest available argument for the labelling discipline described in
the Contributions section: the implementation passed every behavioural
test we had thought to write --- including one specifically checking that anisotropic
embeddings score worse than isotropic ones --- and was still wrong in a way that mattered
more than the bug those tests caught.
